\chapter{谱定理的若干观点}\label{chap:perspectives-spectral}

\section{无限维情形的困难}\label{sec:spectral-infinite-difficulties}

设 \(A\) 是自伴的 \(n\times n\) 矩阵,即 \(A_{kj}=\overline{A_{jk}}\).线性代数中的标准结果断言:存在 \(\mathbb{C}^{n}\) 的标准正交基 \(\{v_{j}\}_{j=1}^{n}\) 以及实数 \(\lambda_{1},\ldots,\lambda_{n}\),使得 \(Av_{j}=\lambda_{j}v_{j}\).(见文献~[24] 第~8 章定理 18 及第~7 章练习 4.)

用与基无关的语言,可以表述同样的结果:设 \(\mathcal{H}\) 是有限维 Hilbert 空间,\(A\) 是其上的自伴线性算子,即对所有 \(\phi,\psi\in\mathcal{H}\),
\[
\langle\phi,A\psi\rangle=\langle A\phi,\psi\rangle.
\]
那么 \(\mathcal{H}\) 存在一个由 \(A\) 的、具有实特征值的特征向量组成的标准正交基.

由于一般 Hilbert 空间也有标准正交基的标准概念,我们可能希望无限维 Hilbert 空间上的自伴算子也有类似结果.不过,简单例子说明,自伴算子可能根本没有特征向量.比如取 \(\mathcal{H}=L^{2}([0,1])\),定义
\begin{equation}
(A\psi)(x)=x\psi(x).
\label{eq:spectral-position-example}
\end{equation}
则 \(A\) 满足 \(\langle\phi,A\psi\rangle=\langle A\phi,\psi\rangle\),但没有特征向量.因为若 \(x\psi(x)=\lambda\psi(x)\),则 \(\psi\) 必须支集在 \(x=\lambda\) 的集合上,而这个集合的测度为零.因此只有 \(L^{2}([0,1])\) 的零元素满足 \(A\psi=\lambda\psi\).

物理学家会说,算子 \(A\) 确实有特征向量,即分布 \(\delta(x-\lambda)\)(见附录~\ref{app:def:distribution}).这些分布满足 \(x\delta(x-\lambda)=\lambda\delta(x-\lambda)\),但不属于 Hilbert 空间 \(L^{2}([0,1])\).这类属于比 \(\mathcal{H}\) 更大空间的“特征向量”称为广义特征向量.虽然广义特征向量并不实际属于 Hilbert 空间,我们仍可能希望在某种意义下它们形成类似标准正交基的对象.第~\ref{sec:momentum-spectral-example} 节将给出这种“基”如何工作的例子.

顺便提及,对于紧自伴算子,这种关于真正特征向量标准正交基的朴素期待确实成立.Hilbert 空间上的算子 \(A\) 称为紧的,是指 \(A\) 对每个有界集的像都有紧闭包;见文献~[34] 第 I 卷定理 VI.16.不过,量子力学中感兴趣的算子并不紧.(当然,即使自伴算子不紧,它仍可能有特征向量标准正交基,例如量子谐振子的 Hamilton 算子,见第~11 章.)

无限维情形中的自伴算子还会带来另一个严重困难.量子力学中的大多数自伴算子都是无界的,即不存在常数 \(C\),使得对所有 \(\psi\) 有 \(\|A\psi\|\le C\|\psi\|\).例如,考虑 \(L^{2}(\mathbb{R})\) 上的位置算子 \(X\),\((X\psi)(x)=x\psi(x)\).若 \(\mathbf{1}_{E}\) 表示集合 \(E\) 的示性函数,则对每个正整数 \(n\),
\[
\|X\mathbf{1}_{[n,n+1]}\|
\ge n\|\mathbf{1}_{[n,n+1]}\|,
\]
所以 \(X\) 不可能有界.再利用闭图定理与第~9.3 节的初等结果可以证明:若 \(A\) 定义在整个 \(\mathcal{H}\) 上,并且对所有 \(\phi,\psi\in\mathcal{H}\) 满足 \(\langle\phi,A\psi\rangle=\langle A\phi,\psi\rangle\),则 \(A\) 必有界.因此,无界自伴算子不可能定义在整个 \(\mathcal{H}\) 上.

于是,我们把 Hilbert 空间 \(\mathcal{H}\) 上的“无界算子”定义为:从 \(\mathcal{H}\) 的稠密子空间——称为 \(A\) 的定义域——到 \(\mathcal{H}\) 的线性算子.此类算子的自伴性比有界情形复杂得多.显然的条件,即对定义域中的所有 \(\phi,\psi\) 要求 \(\langle\phi,A\psi\rangle=\langle A\phi,\psi\rangle\),并不是“正确”条件;具体地说,这个条件不足以保证谱定理适用于 \(A\).相反,对于任意无界算子 \(A\),我们将定义它的伴随 \(A^{*}\),这是一个有自己的定义域的无界算子.无界算子称为自伴,是指 \(A\) 与 \(A^{*}\) 的定义域相同,且它们在共同定义域上相等.也就是说,自伴性不仅要求 \(A\) 与 \(A^{*}\) 在二者都有定义的地方相等,还要求二者的定义域一致.

\section{谱理论的目标}\label{sec:spectral-goals}

在进入谱理论的细节之前,先想一想我们希望谱定理为我们做什么.首先,我们希望谱定理允许我们把各种函数应用到算子上.例如,我们已经看到,可以令 \(\psi(t)=\exp(-it\widehat H/\hbar)\psi_{0}\) 来“求解”含时 Schrödinger 方程.由于 Hamilton 算子 \(\widehat H\) 无界,使用幂级数定义指数并不方便.不过,如果 \(\widehat H\) 有由特征向量 \(\{e_{k}\}\) 构成的真正标准正交基,且相应特征值为 \(\lambda_{k}\),则可以把 \(e^{-it\widehat H/\hbar}\) 定义为满足
\[
e^{-it\widehat H/\hbar}e_{k}
=e^{-it\lambda_{k}/\hbar}e_{k}
\]
的唯一有界算子.

如果 \(\widehat H\) 没有真正的特征向量标准正交基,我们希望谱定理为 \(\widehat H\) 提供一种“函数演算”,即把函数(包括指数函数)应用到 \(\widehat H\) 上的规则.这个函数演算应当具有类似于真正特征向量标准正交基情形中的性质.

其次,我们希望谱定理为测量自伴算子 \(A\) 的结果提供概率分布.回忆 \(A\) 有真正的特征向量标准正交基 \(\{e_{j}\}\) 且特征值为 \(\lambda_{j}\) 时的测量概率.基于例题~\ref{ex:eigenbasis-probability},对 \(\mathbb{R}\) 的任意 Borel 集 \(E\),令 \(V_{E}\) 为 \(A\) 的所有特征值位于 \(E\) 中的特征向量的闭线性包,令 \(P_{E}\) 为投影到 \(V_{E}\) 上的正交投影.那么对任意单位向量 \(\psi\),有
\begin{equation}
\operatorname{Prob}_{\psi}(A\in E)=\langle\psi,P_{E}\psi\rangle.
\label{eq:spectral-measurement-probability}
\end{equation}
特别地,如果特征值两两不同,且 \(\psi=\sum_{j}c_{j}e_{j}\),那么观测到 \(\lambda_{j}\) 的概率是 \(|c_{j}|^{2}\),因为 \(P_{\{\lambda_{j}\}}\) 正是投影到 \(e_{j}\) 上的投影.

如果 \(A\) 没有真正的特征向量标准正交基,我们希望谱定理为每个 Borel 集 \(E\subset\mathbb{R}\) 提供一个投影算子 \(P_{E}\),从而像式~\eqref{eq:spectral-measurement-probability} 那样定义概率.这些投影算子称为\emph{谱投影},相应子空间 \(V_{E}\) 称为\emph{谱子空间}.(于是 \(P_{E}\) 是投影到 \(V_{E}\) 上的正交投影.)直观地说,\(V_{E}\) 可以看成所有特征值属于 \(E\) 的广义特征向量的闭线性包.

在谱定理的第一种形式中,这两个目标都将实现:谱投影由投影值测度给出,函数演算由关于该测度的积分给出.虽然从实用角度看,自伴算子的(广义)特征向量并不是最重要的,但我们将使用直积分的概念提供理解这类特征向量的框架.谱定理的第二种形式把 Hilbert 空间 \(\mathcal{H}\) 分解为关于某个测度 \(\mu\) 的直积分,其中积分对象是自伴算子 \(A\) 的广义特征子空间.对于特定特征值 \(\lambda\),广义特征空间实际上不会是 \(\mathcal{H}\) 的子空间,除非 \(\mu(\{\lambda\})>0\).因此,直积分为“不实际属于 Hilbert 空间的特征向量”提供了严格含义.

\section{阅读指南}\label{sec:spectral-reading-guide}

虽然本书的谱理论部分有些地方不可避免地技术性很强,但它经过设计,读者可以按需要掌握多或少的内容.愿意把一些结果先当作事实接受的读者,可以只看第~\ref{sec:position-spectral-example} 节的位置算子和第~\ref{sec:momentum-spectral-example} 节的动量算子例子,把它们作为谱定理如何工作的原型.希望了解更多细节的读者,可以在第~7 章看到有界算子谱定理的两种表述,并在第~9 章了解无界自伴算子的基本内容.希望完整学习本主题的读者,则可以在第~8 章看到有界算子谱定理两种形式的完整证明,并在第~10 章看到无界算子谱定理的完整证明.

\section{位置算子}\label{sec:position-spectral-example}

作为第一个例子,考虑位置算子 \(X\),
\[
(X\psi)(x)=x\psi(x),
\]
它作用在 Hilbert 空间 \(\mathcal{H}=L^{2}(\mathbb{R})\) 上.与第~\ref{sec:spectral-infinite-difficulties} 节中的类似算子一样,\(X\) 没有真正的特征向量,即没有实际属于 \(\mathcal{H}\) 的特征向量.如果把 \(X\) 的广义特征向量看成 \(\delta(x-\lambda)\)(\(\lambda\in\mathbb{R}\)),那么可以合理猜测:谱子空间 \(V_{E}\) 应由那些“支集在” \(E\) 上的函数组成,即在 \(E\) 的补集上几乎处处为零的函数.(\(\lambda\in E\) 的“函数”\(\delta(x-\lambda)\) 的叠加应是一个支集在 \(E\) 上的函数.)

谱投影 \(P_{E}\) 是投影到 \(V_{E}\) 上的正交投影,可以计算为
\[
P_{E}\psi=\mathbf{1}_{E}\psi,
\]
其中 \(\mathbf{1}_{E}\) 是 \(E\) 的示性函数.于是根据式~\eqref{eq:spectral-measurement-probability},
\[
\operatorname{Prob}_{\psi}(X\in E)
=\langle\psi,P_{E}\psi\rangle
=\int_{E}|\psi(x)|^{2}\,\mathrm{d}x.
\]
这正是我们在第~3 章中关于位置概率分布所预期的公式.

再看 \(X\) 的函数演算.如果 \(f(\lambda)=\lambda^{m}\),则 \(f(X)\) 应是 \(X\) 的 \(m\) 次幂,即乘以 \(x^{m}\) 的算子.因此合理的想法是:对任意函数 \(f\),令 \(f(X)\) 为乘以 \(f(x)\) 的算子.特别地,\(e^{iaX}\) 就是乘以 \(e^{iax}\) 的有界算子.

\section{乘法算子}\label{sec:multiplication-operators}

由于位置算子就是乘以函数 \(x\),所以很容易找出其谱子空间并构造其函数演算.更一般地,若 \(\mathcal{H}=L^{2}(X,\mu)\),且 \(h\) 是 \(X\) 上的实值可测函数,则可在 \(L^{2}(X,\mu)\) 上定义乘法算子 \(M_{h}\):
\[
M_{h}\psi=h\psi.
\]
谱子空间可以取为
\[
V_{E}=\{\psi\mid\psi\text{ 支集在 }h^{-1}(E)\text{ 上}\},
\]
函数演算则定义为
\[
f(A)=\text{乘以 }f\circ h\text{ 的算子}.
\]

谱定理的一种形式可以简单表述如下:可分 Hilbert 空间上的自伴算子 \(A\) 与一个乘法算子酉等价.也就是说,存在某个 \(\sigma\)-有限测度空间 \((X,\mu)\) 和某个可测函数 \(h\),使得 \(A\) 与乘以 \(h\) 的算子酉等价.(见定理~\ref{thm:unbounded-spectral-multiplication}.)虽然这种谱定理表述极其容易说明,但其中涉及直积分的轻微修正有时更好;更多信息见第~7.3 节.

\section{动量算子}\label{sec:momentum-spectral-example}

现在看谱定理如何作用于 \(L^{2}(\mathbb{R})\) 上的动量算子
\[
P=-i\hbar\frac{\mathrm{d}}{\mathrm{d}x}.
\]
\(P\) 的“特征向量”是 \(e^{ikx}\),\(k\in\mathbb{R}\),相应特征值为 \(\hbar k\).虽然 \(e^{ikx}\) 不属于 \(L^{2}(\mathbb{R})\),但 Fourier 变换说明,\(L^{2}(\mathbb{R})\) 中的任意函数都可以表示成这些函数的叠加,即连续版本的线性组合.事实上,Fourier 系数 \(\widehat\psi(k)\) 可以写成 \(\psi\) 与 \(e^{ikx}\) 的“内积”:
\[
\widehat\psi(k)
=(2\pi)^{-1/2}\int_{-\infty}^{\infty}e^{-ikx}\psi(x)\,\mathrm{d}x
=(2\pi)^{-1/2}\langle e^{ikx},\psi\rangle_{L^{2}(\mathbb{R})},
\]
只需忽略 \(e^{ikx}\) 实际不属于 \(L^{2}\) 这一事实.

实际上,物理学家经常断言 \(e^{ikx}/\sqrt{2\pi}\) 构成 \(L^{2}(\mathbb{R})\) 的“连续意义下的标准正交基”.连续意义下的正交归一,意味着在标准正交集定义中用 Dirac \(\delta\) 函数替换通常的 Kronecker delta:
\begin{equation}
\left\langle\frac{e^{ikx}}{\sqrt{2\pi}},\frac{e^{ilx}}{\sqrt{2\pi}}\right\rangle_{L^{2}(\mathbb{R})}
=\delta(k-l).
\label{eq:continuous-orthogonality}
\end{equation}
这里的 \(\delta\) 应满足,对所有连续函数 \(f\),
\[
\int_{-\infty}^{\infty}f(k)\delta(k-l)\,\mathrm{d}k=f(l).
\]
严格地说,\(\delta(k-l)\) 是一个分布(见附录~\ref{app:def:distribution}).

为式~\eqref{eq:continuous-orthogonality} 赋予某种严格含义,虽然 \(e^{ikx}\) 与 \(e^{ilx}\) 的内积没有定义,但可以用如下表达式近似:
\[
\frac{1}{2\pi}\int_{-A}^{A}e^{-ikx}e^{ilx}\,\mathrm{d}x
=\frac{1}{2\pi}\left.\frac{e^{-i(k-l)x}}{-i(k-l)}\right|_{-A}^{A}
=\frac{\sin[A(k-l)]}{\pi A(k-l)}.
\]
可以证明,把上式看成固定 \(A,l\) 时关于 \(k\) 的函数,当 \(A\to\infty\) 时它表现得像 \(\delta(k-l)\).也就是说,对所有足够良好的 \(\psi\),
\begin{equation}
\lim_{A\to\infty}\int_{-\infty}^{\infty}\psi(k)
\frac{\sin[A(k-l)]}{\pi A(k-l)}\,\mathrm{d}k
=\psi(l).
\label{eq:sinc-delta-limit}
\end{equation}

给出式~\eqref{eq:sinc-delta-limit} 的启发式论证.作变量替换 \(k'=k-l\),可把一般问题化为 \(l=0\).再令 \(\kappa=Ak\),所需结果等价于
\begin{equation}
\lim_{A\to+\infty}\int_{-\infty}^{\infty}
\frac{1}{\pi}\frac{\sin\kappa}{\kappa}f\left(\frac{\kappa}{A}\right)\,\mathrm{d}\kappa
=f(0).
\label{eq:sinc-limit-scaled}
\end{equation}
若可以把极限移入积分,则当 \(A\to\infty\) 时 \(f(\kappa/A)\to f(0)\).因此,只需证明
\begin{equation}
\int_{-\infty}^{\infty}\frac{1}{\pi}\frac{\sin\kappa}{\kappa}\,\mathrm{d}\kappa=1.
\label{eq:sinc-integral}
\end{equation}
虽然式~\eqref{eq:sinc-integral} 中的积分不绝对收敛,却是收敛的广义积分;利用围道积分(或查积分表)可得到其值确实为 \(1\).由于式~\eqref{eq:continuous-orthogonality} 本来就是理解 Fourier 变换的启发式方式,我们不再花时间发展上述论证的严格版本.

利用式~\eqref{eq:continuous-orthogonality},至少形式上可以推导 Fourier 变换的许多标准性质,正如利用 \(L^{2}([0,1])\) 中 \(e^{2\pi inx}\) 的正交归一性可以得到 Fourier 级数的性质.更重要的是,Fourier 变换正是把动量算子变为乘法算子的酉变换.为更清晰地看到这一点,作如下简单缩放.

\begin{definition}[动量波函数]\label{def:momentum-wavefunction}
对任意 \(\psi\in L^{2}(\mathbb{R})\),定义
\[
\widetilde\psi(p)=\frac{1}{\sqrt{\hbar}}\widehat\psi\left(\frac{p}{\hbar}\right),
\]
即
\[
\widetilde\psi(p)=\frac{1}{\sqrt{2\pi\hbar}}
\int_{-\infty}^{\infty}e^{-ipx/\hbar}\psi(x)\,\mathrm{d}x.
\]
函数 \(\widetilde\psi(p)\) 称为与 \(\psi\) 关联的动量波函数.
\end{definition}

根据 Plancherel 定理(定理~\ref{app:thm:plancherel-II})和变量替换,如果 \(\psi\) 是单位向量,那么 \(\widehat\psi\) 和 \(\widetilde\psi\) 也都是单位向量.对于任意单位向量 \(\psi\),我们把 \(|\widetilde\psi(p)|^{2}\) 解释为粒子的动量概率密度,正如 \(|\psi(x)|^{2}\) 是粒子位置概率分布.利用命题~\ref{app:prop:fourier-derivatives},可以直接验证,对足够良好的 \(\psi\),有
\begin{equation}
\widetilde{P\psi}(p)=p\widetilde\psi(p).
\label{eq:momentum-multiplication}
\end{equation}
式~\eqref{eq:momentum-multiplication} 意味着酉映射 \(\psi\mapsto\widetilde\psi\) 把动量算子 \(P\) 变成乘以 \(p\) 的算子.也就是说,在“乘法算子”形式的谱定理中,Fourier 变换(按定义~\ref{def:momentum-wavefunction} 的方式缩放)正好实现了谱定理.

用动量波函数,可以像第~\ref{sec:multiplication-operators} 节那样定义 \(P\) 的谱投影和函数演算.对任意 Borel 集 \(E\subset\mathbb{R}\),定义 \(P_{E}\) 为投影到如下函数空间上的正交投影:\(\widetilde\psi(p)\) 在 \(E\) 外几乎处处为零.若 \(f\) 是 \(\mathbb{R}\) 上任意有界可测函数,则可以定义 \(f(P)\):令 \(f(P)\psi\) 为满足
\[
\widetilde{f(P)\psi}(p)=f(p)\widetilde\psi(p)
\]
的 \(L^{2}(\mathbb{R})\) 中唯一元素.
